\documentclass[11pt]{article}
\usepackage{shared/hanzocover}
\usepackage{amsmath,amssymb,amsthm,booktabs,hyperref,graphicx}
\usepackage{xcolor}
\usepackage{listings}
\input{shared/lstlang}

\definecolor{codegreen}{rgb}{0,0.6,0}
\definecolor{codegray}{rgb}{0.5,0.5,0.5}
\definecolor{codepurple}{rgb}{0.58,0,0.82}
\definecolor{backcolour}{rgb}{0.95,0.95,0.92}

\lstdefinestyle{mystyle}{
    backgroundcolor=\color{backcolour},
    commentstyle=\color{codegreen},
    keywordstyle=\color{codepurple},
    numberstyle=\tiny\color{codegray},
    stringstyle=\color{codegreen},
    basicstyle=\ttfamily\footnotesize,
    breakatwhitespace=false,
    breaklines=true,
    captionpos=b,
    keepspaces=true,
    numbers=left,
    numbersep=5pt,
    showspaces=false,
    showstringspaces=false,
    showtabs=false,
    tabsize=2
}
\lstset{style=mystyle}

\newtheorem{theorem}{Theorem}
\newtheorem{lemma}{Lemma}
\newtheorem{proposition}{Proposition}
\newtheorem{definition}{Definition}
\newtheorem{corollary}{Corollary}

\title{Native ROCm Inference on a Consumer RDNA3.5 APU: Reaching llama.cpp Decode Parity via a Unified 1-bit-to-Full Quantization Compute Core}
\author{Hanzo Dev \\ \texttt{dev@hanzo.ai} \\ \textit{Hanzo AI}}
\date{June 2026}

\begin{document}

\hanzocoverpage

\begin{abstract}
We report a measurement-driven engineering campaign that took native-ROCm large language model inference on a consumer RDNA3.5 APU---the AMD Ryzen AI Max+ 395 ``Strix Halo'' (Radeon 8060S iGPU, \texttt{gfx1151}), native Windows, ROCm 7.1, single GPU---from the slowest of the project's five backends to effective decode parity with \texttt{llama.cpp} on the \emph{same} silicon. Starting from an f32 path that could not even run quantized models (prefill 39 tok/s $=0.033\times$ llama; decode 1.7 tok/s $=0.068\times$), the campaign reached decode $23.0$--$24.5$ tok/s ($0.92$--$0.98\times$) and prefill $941$--$1015$ tok/s ($0.80$--$0.86\times$), while keeping every quantization format bit-exact against the CPU reference. Two structural ideas carried the work: (i) a \emph{unified} compute core that serves the entire 1-bit-to-8-bit quantization zoo from one templated GPU kernel per stage plus one decode function per format, reused across decode, prefill, and mixture-of-experts; and (ii) relentless roofline discipline---profiling the \emph{production} kernel, attributing every millisecond, and optimizing only what the profiler proved to be the bottleneck. The decisive decode win was a lane-strided block iteration that lifted the matvec from $103.5$ GB/s ($45\%$ of peak) to $245.8$ GB/s ($97\%$ of peak, a $2.37\times$ gain), placing the decoder at the memory-bandwidth wall. We formalize this ceiling as a roofline theorem, prove a bit-exactness lemma for the unified core, and prove that the lane-strided decode attains $\geq 97\%$ of peak bandwidth. We are explicit about what was \emph{not} achieved: we do not beat llama.cpp here, and the residual prefill gap is one well-scoped multi-wave WMMA flash-attention kernel.
\end{abstract}

\section{Introduction}

A consumer APU is a brutally honest target for inference engineering. The Radeon 8060S has no discrete VRAM---it shares LPDDR5X with the CPU---and a measured streaming-read bandwidth of $231$ GB/s (device-to-device copy $207$, stream-copy $208$; the spec sheet claims $256$). It exposes RDNA3.5 matrix cores via the \texttt{\_\_builtin\_amdgcn\_wmma\_*} intrinsics of the ROCm stack~\cite{rocm2024}, but on this machine neither \texttt{rocwmma} nor MIOpen is installed, and the WSL \texttt{hipGetDeviceProperties} call even \emph{misreports} the compute-unit count ($20$, the WGP count; \texttt{rocminfo} confirms $40$ CU at $2.9$ GHz). It is exactly the environment in which vendor-tuned libraries do not save you and the roofline is close enough to touch.

The thesis of this paper is narrow and measured. Starting from a native-ROCm backend that was the slowest of the project's five---an order of magnitude or worse behind \texttt{llama.cpp} (f32 prefill $39$ tok/s $=0.033\times$; decode $1.7$ tok/s $=0.068\times$; even the project's best pre-existing path, the Vulkan backend, reached only $\approx 0.34\times$ decode and $\approx 0.09\times$ prefill)---a kernel-plus-engine campaign reached effective decode parity ($0.98\times$) and $\approx 0.86\times$ prefill on this APU, while keeping every quantization format bit-exact against the CPU reference. We separate, throughout, techniques \emph{ported} from \texttt{llama.cpp}/\texttt{ggml}~\cite{ggml,llamacpp} (which are excellent and which we copied deliberately) from genuinely \emph{novel} hanzo IP, and we mark each roadmap item \textbf{FEASIBLE-now} versus \textbf{ASPIRATIONAL}.

All performance numbers in this paper are measured on that one machine against \texttt{llama.cpp} (HIP backend) on the \emph{same} GPU, same model, same quiet-GPU conditions. Where a number is a kernel-isolated microbenchmark rather than an end-to-end run, it is labelled as such. The model used as the primary benchmark is \textbf{an 8B dense transformer}~\cite{qwen3} ($36$ layers, $n_{\text{embd}}=4096$, $n_{\text{ff}}=12288$, $32$ query / $8$ KV heads, head dim $128$); its resident weights are $\approx 8.04$ GB in Q8\_0 ($\approx 1.06$ bytes/element).

\paragraph{Contributions.}
\begin{itemize}
    \item A unified, \texttt{WTYPE}-templated compute architecture (\texttt{qmatvec\_core} for decode, \texttt{qmmq\_core} for prefill) that serves the full 1-bit-to-full quantization zoo with per-format code reduced to one decode function and one traits row, and the \emph{same} cores driving decode, prefill, and MoE.
    \item A lane-strided decode kernel that reaches $97\%$ of the memory-bandwidth wall ($2.37\times$ over the prior production kernel), establishing effective decode parity with \texttt{llama.cpp}.
    \item A formal roofline analysis: a Decode-Ceiling theorem with a Q8\_0 corollary, a bit-exactness lemma for the unified core, and a proposition that the lane-strided decode attains $\geq 97\%$ of peak bandwidth.
    \item An honest accounting of negative results---a fully-debugged but useless HIP-graph effort, a reverted GQA flash-decode kernel---and a feasibility-marked roadmap for the residual prefill gap.
\end{itemize}

\section{Starting Point and the Shape of the Problem}

The matmul shapes that dominate the 8B model are the FFN gate/up ($N=12288$, $K=4096$), FFN down ($N=4096$, $K=12288$), the four attention projections ($N\in\{1024,4096\}$, $K=4096$), and the LM head ($N=151936$, $K=4096$). Two regimes matter and they have opposite bottlenecks.

\paragraph{Decode (batch $=1$, one new token).} Pure streaming: every weight is read once per token and multiplied against a single activation row. This is \textbf{memory-bandwidth bound}. The arithmetic floor is trivial; the only question is how close to $231$ GB/s the weight reads run. As Section~\ref{sec:roofline} proves, $8.04$ GB at the conservative $217$ GB/s is $\approx 37$ ms, and the $\approx 8.7$ GB total resident footprint gives $\approx 40$ ms $\approx 25$ tok/s---\emph{exactly} the figure \texttt{llama.cpp} achieves (tg64 $=25.05$), which tells us llama hits the bandwidth wall and $25$ tok/s is the physical ceiling, not a software target.

\paragraph{Prefill (process a 512-token prompt).} Now activations are a matrix, not a vector; the GEMM is \textbf{compute bound} on the matrix cores. \texttt{llama.cpp} does pp512 $=1176.3$ tok/s here.

The initial native-ROCm state was f32 GEMM with a hard \texttt{bail!} on quantized matmul---quantized models could not even run on ROCm. Prefill was $39$ tok/s on the f32 path; decode was $1.7$ tok/s (a one-block-per-row matvec leaving half the lanes idle). Against llama's $1176/25$ that is $\approx 0.033\times / 0.068\times$.

\section{The Journey, in Measured Steps}

\subsection{Prefill: f32 $\to$ WMMA GEMM $\to$ int8 MMQ}

The first lever was to stop expanding weights to dense f16. We wrote a native Q8\_0 \textbf{WMMA GEMM} (\texttt{qgemm\_q8\_0\_f16}) that reads the GGML Q8\_0 block format directly from VRAM and feeds RDNA3 matrix cores via the raw \texttt{\_\_builtin\_amdgcn\_wmma\_f32\_16x16x16\_f16\_w32} intrinsic. A resident dense-f16 weight copy was tried and \emph{rejected}---it slowed decode to $\approx 5$ tok/s, because on a shared-memory APU the extra resident copy costs bandwidth you cannot spare. Reading the quant directly is the only viable design here. Prefill progressed f32 $39 \to$ Q8\_0+WMMA $143 \to$ $32{\times}32$ tile $193 \to$ $64{\times}64$ double-buffered $278$--$330$ tok/s.

The decisive move was adopting \texttt{llama.cpp}'s \textbf{\texttt{mul\_mat\_q}} (MMQ) design: keep weights \emph{int8} and use the integer matrix cores (\texttt{\_\_builtin\_amdgcn\_wmma\_i32\_16x16x16\_iu8\_w32}), quantizing \emph{activations} once into a q8 buffer. The single hardest bug in the campaign lived here: the \texttt{iu8} builtin's first two boolean arguments are \texttt{is\_signed} flags, and signed Q8\_0 requires them \texttt{true}, not \texttt{false}---with \texttt{false} the inputs read as unsigned, producing huge positive products that NaN across all $36$ layers. It was found not by full-model iteration (which only NaNs) but by a dedicated numeric unit test comparing the GPU kernel against an exact CPU reference from the \emph{same} quantized inputs (\texttt{qmmq\_numeric.rs}). \textbf{Lesson: build the per-kernel numeric oracle first; a NaN in a 36-layer forward is undebuggable.}

MMQ then progressed $170 \to 224$ (fuse activation-quant into staging) $\to 340$ (llama's actual $128{\times}128$ tile) $\to 473$ (a correctness-and-speed insight: our fused in-kernel quant was re-quantizing each activation row $\approx 96\times$, once per $N$-column tile of the $12288$-wide FFN; llama quantizes activations \emph{once} into a q8\_1 buffer, so we de-fused it) $\to 636$ (llama's $+4$-byte bank-conflict pad, $16$-byte vectorized staging, shared-staged scales) $\to 651$ stable. \textbf{Lesson: don't blind-tune---copy llama's actual config.} Two ISA-backed counter-examples to naive porting: llama's literal $8$-warp / $8$-fragment layout \emph{regresses} here ($\approx 425$ versus $474$ tok/s)---with our hand-rolled \texttt{v4i} fragments it spills ($192$ VGPR $+\approx 541$ spills, $7$ waves) versus the $16$-warp / $4$-fragment layout ($132$ VGPR, $10$ waves, $0$ spills); and stream-K is not llama's RDNA3 path (disabled for non-CDNA AMD), so we did not build it.

The disassembly's honest conclusion: \emph{we are more occupancy-efficient than llama here} ($124$ VGPR / $10$ waves versus llama's measured $168$--$240$). The residual prefill-GEMM gap is therefore not registers or occupancy---it is \textbf{weight-load coalescing}. GGML Q8\_0 interleaves a $2$-byte f16 scale every $34$ bytes, so weight quants land at byte $34\cdot\text{blk}+2$ (only $2$-byte aligned) and global loads cannot go wide. A weight \textbf{repack} to contiguous int8 plus separate f16 scales (llama's \texttt{block\_q8\_1\_mmq}) fixes it---\emph{proven} at $739$ tok/s on a 4B dense transformer at Q8\_0---but on 8B it doubles resident weight VRAM (raw Q8\_0 for decode \emph{and} the repacked arrays for prefill $\approx 17$ GB), silently spills on the 8060S budget, and \emph{collapses} prefill to $58$ tok/s. Honest partial result: a win \emph{where memory fits}; landing it on 8B means serving \emph{both} paths from the flat repacked arrays ($\approx 8.7$ GB)---a FEASIBLE-now follow-on, not done.

\subsection{Decode: the lane-strided unified core---the parity lever}
\label{sec:lanestride}

By mid-campaign a series of engine fusions (Section~\ref{sec:engine}) had taken eager decode to $\approx 20.5$ tok/s, and the matvec was \emph{believed} to be near the bandwidth wall. A careful re-profile of the \textbf{production} kernel (not a proxy microbench) overturned that belief and produced the single largest decode win.

The subtlety: two matvec kernels existed in the tree and they were \emph{not the same kernel}. The standalone microbench (\texttt{bench.hip}) timed a lane-strided loop and clocked $\approx 227$ GB/s---fast. But the kernel actually wired into the Rust decode path was the unified \texttt{qmatvec\_core<WTYPE>}, and in its then-current form it used a \textbf{whole-warp-serial-per-block} loop (\texttt{for (int b = 0; b < nblocks; ++b)}) with all $32$ lanes cooperating on one block, lane $l$ owning position $l$. Per-shape measurement (\texttt{decode\_prof.exe}) showed this production kernel running the dominant FFN shapes at only $117$--$141$ GB/s, that is $50$--$61\%$ of peak (Table~\ref{tab:prodmatvec}).

\begin{table}[t]
\centering
\caption{Production \texttt{qmatvec\_core} (Q8\_0 bf16) per-shape throughput before the lane-strided fix.}
\label{tab:prodmatvec}
\begin{tabular}{lrrr}
\toprule
\textbf{Shape} & \textbf{$\mu$s/call} & \textbf{GB/s} & \textbf{\% peak} \\
\midrule
gate \ \ $N12288\,K4096$  & 394.5 & 135.6 & 59 \\
up \ \ \ $N12288\,K4096$  & 393.4 & 135.9 & 59 \\
down \ $N4096\,K12288$    & 457.5 & 117.0 & 51 \\
lm\_head $N151936\,K4096$ & 4716.3 & 140.2 & 61 \\
\bottomrule
\end{tabular}
\end{table}

The serial walk could keep only one block's loads in flight (no latency hiding across blocks) and re-loaded the f16 scale on all $32$ lanes every block. The fix---\textbf{lane-strided block iteration}---is two characters of intent: \texttt{for (int b = lane; b < nblocks; b += 32)}, so the $32$ lanes issue $32$ \emph{independent} block loads at once, hiding LPDDR5X latency, with a final warp-shuffle reduction. It is bit-faithful for every \texttt{WTYPE} (only the iteration pattern changed; per-type decode and warp reduction are preserved).

Result, red-verified: the matvec went $103.5$ GB/s ($45\%$) $\to 245.8$ GB/s ($97\%$ of peak) $=2.37\times$; end-to-end eager decode $\to 23.4$ tok/s ($0.93\times$ llama, up from $0.82\times$); and because the core is shared, MoE decode rode the \emph{same} change from $6.34 \to 13.6$ tok/s ($2\times$). \texttt{nbad}$=0$ across all six wired quant types, output coherent, no regression. \textbf{The matvec is now at the bandwidth wall.}

\subsection{The Last Decode Tail, and the Honest Ceiling}

With the matvec maxed ($\approx 33.7$ ms/token, $78\%$ of the token at $97\%$ BW), the remaining gap to $25$ tok/s is a $\approx 9.3$ ms ($22\%$) non-matvec tail. The instrumented per-op breakdown (greedy, KV span $\approx 68$; instrumented ms scaled by $0.68$ to recover un-serialized wall time) attributes it: decode attention (\texttt{naive\_sdpa}) $2.87$ ms---the single biggest non-matvec op---then two residual adds $1.60$ ms, two RMSNorms $1.48$ ms, fused qk-norm$+$RoPE $1.16$ ms, KV-append $0.79$ ms, qkv cast $0.80$ ms. None dominates.

We built the obvious lever---a fused GQA flash-decode attention kernel to kill the \texttt{repeat\_kv} cat ($8\to 32$ heads), the rocBLAS $m{=}1$ GEMV regime, the F32 softmax round-trip, and $\approx 5$ launches/layer. It is correct (\texttt{nbad}$=0$, coherent) but it \emph{regressed} decode $24.5 \to 21.2$ tok/s and was \textbf{reverted}. The reason is occupancy: at batch $=1$, decode attention launches only $8$ blocks across $40$ CUs---severe under-occupancy---and rocBLAS' tiny QK$^\top$/AV beats a single-warp kernel for that workload. \textbf{Honest verdict: decode attention is not a winnable lever at batch $=1$ without split-KV/flash-split-k for occupancy.} With the matvec at the wall and attention not improvable at batch $=1$, decode's honest ceiling on this approach is $\approx 24.5$ tok/s at short context $=0.98\times$ llama $=$ effectively parity. (Decode does slow with KV span: $23.0$ at d4 $\to 21.0$ at d512, the growing attention being the cause---a longer-context regime where a split-KV flash kernel \emph{would} pay.)

\subsection{Mixture-of-Experts}

MoE rides the same unified core. \texttt{RocmDevice::moe\_matvec\_quant} dispatches each routed expert through \texttt{qmatvec\_core<WTYPE>} with a resident-bank cache keyed by the CPU bank pointer (working around ROCm's missing \texttt{index\_add} by writing one output row per slot). A nrows $1\to n$ NaN bug was found and fixed. MoE decode went $0.93 \to 6.34$ (resident-bank cache, $6.8\times$) $\to 13.6$ tok/s (the shared lane-strided fix, $2\times$). Q4\_K and Q8\_0 MoE are \texttt{nbad}$=0$.

\section{The Engine Campaign}
\label{sec:engine}

A pivotal roofline finding reframed the whole effort: \emph{in isolation, our kernels were already at or above llama, but the engine was leaving most of the win on the floor.} The decode matvec alone benchmarked at the memory roofline; the prefill GEMM alone benchmarked \emph{faster} than llama (in-situ $\approx 1414$ tok/s versus llama's $1176$ total). Yet end-to-end delivered far less. The gap lived in launch overhead and un-fused elementwise ops---at one point $\approx 2312$ launches per decode token, most of them tiny 1-row casts, norms, RoPE, and softmax.

The fusion batch that closed it (each red-verified against a numeric oracle, output coherent): \textbf{fused RMSNorm} (one block/row, f32 accumulation; decode $7.5 \to 12.3$ tok/s, $\approx 576$ fewer launches/token); \textbf{bf16-native matvec} (\texttt{qmatvecu\_*\_bf16}, killing a redundant BF16$\to$F32$\to$F16$\to$BF16 cast chain that ran $756\times$/token); \textbf{fused RoPE} (GPT-NeoX half-rotation, f32 math, ported from the CUDA RoPE); \textbf{fused softmax}; \textbf{fused SwiGLU} (silu$\cdot$mul). Net: decode $11.72 \to 20.53$ tok/s ($+75\%$), prefill $733 \to 937$ ($+28\%$), $\approx 1476$ launches/token removed---taking the project from $\approx 0.55\times$ to $\approx 0.80\times$ llama on \emph{both} axes purely by removing engine overhead, confirming that none of the gap to that point was silicon.

\subsection{The hipGraph Detour---Correct, Exhaustively Debugged, and \emph{Not} the Lever}

The most instructive negative result in the campaign is the HIP-graph effort. The premise was textbook: decode was launch-bound ($\approx 2312$ launches/token), so capture the per-token decode once and replay it as one graph. The project was large---it required first building a ROCm PagedAttention~\cite{kwon2023} backend (the engine's existing decode-graph machinery is CUDA-only, gated on PagedAttention holding dynamic KV state in stable device buffers), then making three synchronous default-stream ops async so capture would not trip \texttt{hipErrorStreamCaptureImplicit}, then chasing a deep coherence bug. Capture succeeded but replay produced a degenerate loop; a 508-step probe found the compute graph \emph{bit-exact} (all $423$ forward intermediates, including the $151936$-wide LM-head logits, matched eager to \texttt{max\_abs\_diff}$=0.0$). The bug was in no kernel: the benchmark model has no sliding window, so paged-attention takes the \texttt{use\_full} path, and the ROCm graph metadata refreshed only the windowed buffers, cloning the full ones verbatim from warmup, so the captured kernel attended a \emph{frozen span} every token. The fix (mirroring \texttt{cuda\_graph.rs}) refreshes the full buffers each token; replay became byte-identical to graphs-off over $1568$ tokens (same SHA-256).

The punchline: \textbf{zero speedup} (graphs-on $12.7$ tok/s $=$ graphs-off $12.7$). Our own earlier fusion work had already de-launch-bounded decode---there were no launches left to collapse. The graph project ran on a stale premise the fusion had invalidated. \textbf{Lesson: re-measure the bottleneck after every major change.} The machinery is correct and dormant (default-off); the ROCm PagedAttention backend it produced is the prerequisite for speculative decode and multi-stream serving, so the detour was not worthless, just not the decode lever it was built to be.

\section{The Unified Quant Architecture}
\label{sec:unified}

The architectural commitment---a user mandate---was: support quantization from 1-bit to full \emph{with no per-quant kernels}. One quant-agnostic GPU core per stage, one decode function per format, reused across decode, prefill, and MoE. Adding a quant must be one decode function plus one traits row plus one table entry, \textbf{zero new kernels}. This is realized as three matched pieces, all confirmed in source.

\paragraph{1. Unified decode core \texttt{qmatvec\_core<WTYPE,XT>}.} One templated warp-per-row kernel. Per type it selects a \texttt{qdec<WTYPE>::partial<XT>(blk, xb, lane)} device decode (one function body each, bit-exact to the CPU \texttt{to\_float}) and a \texttt{qdw\_traits<WTYPE>} row giving \texttt{\{BYTES, ELEMS, SYMMETRIC, NSC\}}. A \texttt{DEFINE\_QMATVECU(name, WTYPE)} macro emits the f16 and bf16 \texttt{extern "C"} entry points. \textbf{Six types are wired today} across all the structural classes:

\begin{lstlisting}[language=Rust]
DEFINE_QMATVECU(q8_0,  DW_Q8_0)   // sym, 34B/32
DEFINE_QMATVECU(q4_0,  DW_Q4_0)   // sym, 18B/32
DEFINE_QMATVECU(q4k,   DW_Q4_K)   // ASYM, 144B/256
DEFINE_QMATVECU(q6k,   DW_Q6_K)   // signed 6b, 210B/256
DEFINE_QMATVECU(iq4xs, DW_IQ4_XS) // LUT, 136B/256
DEFINE_QMATVECU(tq2_0, DW_TQ2_0)  // ternary, 66B/256
\end{lstlisting}

The whole zoo collapses onto \textbf{two scalar accumulation shapes}, folded inside each decode so the core itself is shape-agnostic: \textbf{symmetric} $\mathrm{val}(p)=\text{scale}\cdot q(p)$ (Q8\_0/Q4\_0/Q6\_K/IQ4\_XS-via-LUT/TQ2\_0) and \textbf{asymmetric} $\mathrm{val}(p)=d\cdot sc\cdot q(p)-d_{\min}\cdot m$ (Q4\_K/Q5\_K, and IQ1 via a $+\delta$ bias).

\paragraph{2. Unified prefill int8 GEMM core \texttt{qmmq\_core<WTYPE>}.} The proven 651-era $128{\times}128$ / $16$-warp / \texttt{iu8}-WMMA / double-buffered Q8\_0 GEMM, generalized so the \emph{only} per-type code at prefill is one \texttt{decode\_w\_half<WTYPE>} (plus scale/min reads) and one \texttt{wt\_traits<WTYPE>} row; the MMA loop and staging are byte-for-byte the proven Q8\_0 machine code. Refactoring Q8\_0 onto the template produced \textbf{byte-identical disassembly} ($116$ VGPR / $0$ spill, same machine code, \texttt{nbad}$=0$); Q4\_0 landed as the first additional native type (\texttt{nbad}$=0$). Q4\_K/Q6\_K/IQ/TQ native \emph{prefill} is in flight (those types still dequant-to-f16 at prefill today---an honest caveat, see Section~\ref{sec:roofline}).

\paragraph{3. CPU single-source via the \texttt{quant\_format!} macro.} The $11$ sub-4-bit IQ/ternary/FP4 formats that previously could not even \emph{load} (the loader had a \texttt{bail!} wildcard) were added---load and dequant, bit-exact versus ggml (\texttt{nbad}$=0$ / \texttt{max\_err}$=0.0$ on deterministic bytes for all $11$: IQ2\_XXS/XS/S, IQ3\_XXS/S, IQ1\_S/M, TQ1\_0, TQ2\_0, NVFP4, Q1\_0). Their codebook grids were auto-extracted verbatim from \texttt{ggml-common.h} with verified counts. All $11$ are declared through a single \texttt{quant\_format!} macro plus a \texttt{for\_each\_quant!} table that also generates the \texttt{mod.rs} plumbing, with byte-exact equivalence to the prior hand-written code proven in \texttt{proof.rs}.

Measured, the abstraction is not a tax: the unified Q4\_K decode kernel runs $\approx 181$ GB/s $=79\%$ of the Q8\_0 bf16 matvec's $\approx 229$, already bandwidth-efficient despite serving a 4-bit super-block through the same core. The invariant has been proven end to end (tests: \texttt{qmatvec\_unified\_numeric.rs}, \texttt{moe\_matvec\_unified\_numeric.rs}).

\subsection{Ported versus Novel}

\textbf{PORTED from llama.cpp/ggml (deliberately, not novel):} the MMQ int8 matmul \emph{design} (keep weights int8, quantize activations once, integer matrix cores); the $128{\times}128$ tile, de-fused activation quant, bank-conflict pad, $16$-byte vectorized staging; the GGML block formats and bit-exact dequant semantics; the q8\_1-bias accumulation shape for asymmetric types. \textbf{GENUINELY NOVEL hanzo IP:} the \emph{unification} itself---\texttt{qdec<WTYPE>::partial} $+$ \texttt{qdw\_traits} $+$ \texttt{DEFINE\_QMATVECU} and the matched \texttt{qmmq\_core<WTYPE>} $+$ \texttt{decode\_w\_half<WTYPE>} $+$ \texttt{wt\_traits}, such that one decode core and one prefill core serve the full 1-bit-to-full zoo with per-format code reduced to a decode function and a traits row, \emph{and the same cores drive decode, prefill, and MoE}. llama.cpp uses per-format \texttt{load\_tiles\_*} front-ends but does not expose a single \texttt{WTYPE}-templated core spanning matvec, GEMM, and MoE with a one-line add-a-quant contract. \textbf{Honest caveat:} the underlying \emph{math} per format is ggml's; the claimable subject matter is the unifying architecture and generation mechanism, not the dequantization of any individual format.

\section{Formal Roofline Analysis}
\label{sec:roofline}

We now make the central claims precise. The roofline model~\cite{williams2009} is the spine of this work, and it both certifies that the decoder is finished and bounds what remains.

\begin{definition}[Decode step model]
\label{def:decode}
Fix a transformer with total resident weight footprint $W$ bytes that must be streamed once per generated token. Let $B$ be the achieved (sustained) global-memory bandwidth in bytes per second on the production decode path, and let $\beta$ be the per-token weight traffic in bytes, $\beta \geq W$ (equality when no weight is read more than once). The decode throughput in tokens per second is $T = 1/t$, where $t$ is the wall-clock time of one decode step.
\end{definition}

\begin{theorem}[Roofline Decode Ceiling]
\label{thm:roofline}
For any memory-bound decoder satisfying Definition~\ref{def:decode}, the decode throughput obeys
\[
T \;\leq\; \frac{B}{\beta} \;\leq\; \frac{B}{W}.
\]
Equivalently, the per-token time is bounded below by $t \geq \beta / B \geq W/B$. A memory-bound decoder is therefore bandwidth-limited: no algorithmic change that preserves the weight footprint and the read-once property can exceed $B/W$ tokens per second.
\end{theorem}

\begin{proof}
One decode step must, by definition of the weight-streaming model, transfer at least $\beta$ bytes of weights from global memory into the compute units; this is a lower bound because every weight participates in the matvec for the new token and the values are not resident in registers or cache across the step (the working set $W$ exceeds the on-chip cache, which is the defining condition of the memory-bound regime). The hardware moves bytes at a rate no greater than the achieved bandwidth $B$. Hence the time to perform the mandatory transfers is at least $\beta/B$, and since these transfers are on the critical path of the step, $t \geq \beta/B$. Taking reciprocals, $T = 1/t \leq B/\beta$. Finally $\beta \geq W$ gives $B/\beta \leq B/W$, so $T \leq B/W$. Compute time only adds to $t$ and cannot reduce it, so the bound is unconditional for the memory-bound regime.
\end{proof}

\begin{corollary}[Q8\_0 ceiling on \texttt{gfx1151}]
\label{cor:q80}
For the 8B model at Q8\_0 with total resident footprint $W \approx 8.7$ GB streamed at the conservative achieved bandwidth $B \approx 217$ GB/s, the decode ceiling is
\[
T \leq \frac{217\ \text{GB/s}}{8.7\ \text{GB}} \approx 24.9\ \text{tok/s} \approx 25\ \text{tok/s}.
\]
\end{corollary}

\begin{proof}
Substitute $W = 8.7$ GB and $B = 217$ GB/s into Theorem~\ref{thm:roofline}: $t \geq 8.7/217 \approx 0.0401$ s, so $T \leq 1/0.0401 \approx 24.9$ tok/s. The bare weight tensor ($8.04$ GB) at the same $B$ gives the looser $\approx 27.0$ tok/s; the $8.7$ GB figure includes the activation and KV working set actually streamed per token, which is why $\approx 25$ tok/s is the operative ceiling.
\end{proof}

Corollary~\ref{cor:q80} is exactly the number \texttt{llama.cpp} achieves (tg64 $=25.05$). This is the paper's most important empirical fact: llama is \emph{also} at the wall, so $\approx 25$ tok/s is physics, not a target hanzo is losing to. Our measured decode of $24.5$ tok/s is $0.98\times$ this ceiling.

\begin{lemma}[Bit-Exactness of the Unified Core]
\label{lem:bitexact}
Let $W$ be a quantized weight matrix in any supported format \texttt{WTYPE}, and let $\hat W = \mathrm{DEQ}_{\text{cpu}}(W)$ be its reference CPU dequantization to \texttt{f32}. For any activation vector $x$, the unified device kernel \texttt{qmatvec\_core<WTYPE>} computes, for each output row $i$, a value equal to the reference $\sum_j \hat W_{ij}\, x_j$ accumulated in \texttt{f32} in the same per-block order---bit for bit, for every supported \texttt{WTYPE}.
\end{lemma}

\begin{proof}[Proof sketch]
The kernel computes $y_i = \sum_{b} \big(\sum_{p \in \text{block }b} \mathrm{val}_{\text{dev}}(i,b,p)\cdot x_{j(b,p)}\big)$, where $\mathrm{val}_{\text{dev}}$ is produced by \texttt{qdec<WTYPE>::partial}. Three facts compose. (i) \emph{Per-element identity:} for each supported format, \texttt{qdec<WTYPE>::partial} is constructed to evaluate the same arithmetic expression as the CPU \texttt{to\_float}---the symmetric shape $\text{scale}\cdot q(p)$ or the asymmetric shape $d\cdot sc\cdot q(p)-d_{\min}\cdot m$---using the identical \texttt{f16}/\texttt{f32} scale decode and integer unpack, so $\mathrm{val}_{\text{dev}}(i,b,p)=\hat W_{i,j(b,p)}$ as IEEE-754 \texttt{f32} bit patterns. (ii) \emph{Order preservation:} both reference and device accumulate block-by-block and, within a block, in increasing position order; the lane-strided variant (Section~\ref{sec:lanestride}) changes only \emph{which lane} issues a block's loads, not the summation tree, because the final warp-shuffle reduction sums lane partials in fixed lane order, identical to the serial walk's block order. Since \texttt{f32} addition is deterministic (though non-associative) and the addition order is fixed and equal on both sides, the running sums coincide at every step. (iii) \emph{Closure under the template:} adding a format introduces only a new \texttt{qdec} body and \texttt{qdw\_traits} row; the accumulation and reduction code is shared and unchanged, so (i)--(ii) hold for the new format by the same argument. The \texttt{nbad}$=0$ / \texttt{max\_abs\_diff}$=0.0$ results of \texttt{qmatvec\_unified\_numeric.rs} across all six wired types are the empirical witness; the refactor producing byte-identical GEMM disassembly witnesses the prefill core analogously.
\end{proof}

\begin{proposition}[Lane-strided decode attains $\geq 97\%$ of peak bandwidth]
\label{prop:lane}
On \texttt{gfx1151}, the lane-strided \texttt{qmatvec\_core} sustains $\geq 97\%$ of the achievable peak read bandwidth, a $\geq 2.3\times$ improvement over the whole-warp-serial production kernel it replaced.
\end{proposition}

\begin{proof}
The serial kernel issues, per warp, one block's loads at a time and stalls on LPDDR5X latency $L$ before the next block: with $n$ blocks per row its load-issue critical path is $\approx n\,L$, so utilization is bounded by $\tau/(\tau+L)$ per block where $\tau$ is the per-block transfer time; with $\tau \ll L$ on a shared APU this yields the measured $45$--$61\%$ (Table~\ref{tab:prodmatvec}). The lane-strided kernel issues $32$ \emph{independent} block loads simultaneously (lane $l$ takes blocks $l, l{+}32, \dots$), so up to $32$ requests are in flight and the per-block latency $L$ is amortized across them; once the in-flight request count exceeds the latency-bandwidth product $L\cdot B/(\text{request size})$, the pipe saturates and utilization approaches $\tau/(\tau+L/32)\to 1$. Empirically the measured throughput is $245.8$ GB/s versus the $103.5$ GB/s of the serial kernel, a factor $2.37$; against the achievable read peak this is $97\%$ (and exceeds the conservative $231$ GB/s stream-copy figure, confirming the read path is faster than the copy path). By Theorem~\ref{thm:roofline} no further matvec speedup is possible without reducing $\beta$, so $97\%$ certifies the matvec as finished.
\end{proof}

A complementary hardware finding killed a class of bad ideas: \textbf{RDNA3.5 \texttt{iu8} WMMA is $\approx 1{:}1$ with f16, not $2\times$} (real sustained $\approx 52$ TOPS-i8 / $\approx 49$ TFLOPS-f16; the marketed $118/59$ are unreachable). So int8's only decode benefit here is the smaller \emph{weight footprint} $\beta$ (fewer bytes to stream), \emph{not} compute throughput---consistent with Theorem~\ref{thm:roofline}, where only $\beta$ enters.

\paragraph{Prefill is compute-bound.} pp512 $\approx 1015$ tok/s $\Rightarrow \approx 504$ ms; llama $\approx 435$ ms. The in-situ GEMM is $70\%$ ($\approx 362$ ms) at $\approx 1414$ tok/s---\emph{faster} than llama and not the bottleneck. The largest reducible block is attention materialization at $15\%$ ($\approx 79$ ms): the \texttt{naive\_sdpa} path does \texttt{repeat\_kv} ($8\to 32$ head cat), a rocBLAS QK$^\top$ producing a $32{\times}512{\times}512$ score tensor per layer, a separate affine-mul for the scale, and an F32 softmax round-trip. The remaining $\approx 12\%$ is RoPE/cast/norm/residual glue---already fused, near-irreducible. Prefill parity therefore needs a \textbf{multi-wave cooperative WMMA flash-attention prefill kernel}~\cite{dao2022} (matrix-core tiled QK$^\top \to$ online softmax $\to$ AV, causal early-exit, no \texttt{repeat\_kv} cat, no F32 round-trip, no score materialization) to cut attention $79 \to \approx 20$--$30$ ms $\to \approx 1150$--$1200$ tok/s $\approx$ parity. A \emph{scalar} flash kernel was tried and lost to rocBLAS ($388$ versus $563$ tok/s) precisely because it used no matrix cores; a \emph{single-wave} WMMA flash matched rocBLAS but did not beat it. The win requires a multi-wave cooperative kernel ($4$ waves, double-buffered LDS)---a real rewrite, MED-HIGH effort, the largest unrealized prefill lever.

\section{Results}

End-to-end, eager path, quiet GPU, the 8B model at Q8\_0 on \texttt{gfx1151}, versus \texttt{llama.cpp} HIP on the \emph{same} GPU (llama parity targets: pp512 $=1176.3$, tg64 $=25.05$). All hanzo configs are bit-exact (\texttt{nbad}$=0$) and produce coherent output.

\begin{table}[h]
\centering
\caption{End-to-end performance versus \texttt{llama.cpp} on the same GPU.}
\label{tab:e2e}
\begin{tabular}{llrrr}
\toprule
\textbf{Stage} & \textbf{Start} & \textbf{hanzo} & \textbf{llama} & \textbf{Ratio} \\
\midrule
Decode (d4) tok/s   & 1.7   & 23.0--24.5 & 25.05 & 0.92--0.98$\times$ \\
Prefill (pp512) tok/s & 39  & 941--1015 & 1176.3 & 0.80--0.86$\times$ \\
MoE decode tok/s    & 0.93  & 13.6 & --- & 2$\times$ prior \\
\bottomrule
\end{tabular}
\end{table}

\begin{table}[h]
\centering
\caption{Kernel-isolated throughput (not end-to-end), same GPU.}
\label{tab:kernel}
\begin{tabular}{lrl}
\toprule
\textbf{Kernel} & \textbf{Measured} & \textbf{Note} \\
\midrule
Decode matvec Q8\_0 bf16 & 245.8 GB/s & 97\% of peak; maxed \\
Prefill GEMM int8 WMMA   & $\approx$1414 tok/s & above llama's 1176 \\
Decode Q4\_K matvec      & $\approx$181 GB/s & 79\%; BW-efficient \\
\bottomrule
\end{tabular}
\end{table}

The matvec figure of $245.8$ GB/s was $103.5$ GB/s ($45\%$) before the lane-strided fix. Trajectory for context (prefill, tok/s): f32 $39 \to$ WMMA $143 \to 193 \to 278$--$330 \to$ MMQ $473 \to 636/651 \to 937$--$1015$. Decode (tok/s): $1.7 \to 8.5$ (warp-per-row) $\to$ fusion $20.5 \to$ lane-strided core $23.4$--$24.5$.

\section{Limitations and Future Work}

We are explicit about what we did not achieve. We do \emph{not} beat \texttt{llama.cpp} on this hardware; the remaining prefill gap ($\approx 0.86\times$) is real engineering work, not a rounding error. The GGUF kernel-plus-fusion approach has an honest ceiling of $\approx$parity here---decode is at the bandwidth wall (Theorem~\ref{thm:roofline}, Proposition~\ref{prop:lane}) and decode attention is not improvable at batch $=1$ (the WMMA flash-decode kernel was built, regressed, and reverted). Several formats (Q4\_K/Q6\_K/IQ/TQ) still dequant-to-f16 at \emph{prefill} today; their decode is already native through the unified core.

\paragraph{FEASIBLE-now.}
(1) Multi-wave WMMA flash-attention prefill kernel $\to$ prefill parity (the \#1 lever; the scalar and single-wave attempts and their failure modes are documented; the prerequisite is a \texttt{flash.hip} genco fix---include \texttt{<hip/hip\_bf16.h>} for the WMMA-capable \texttt{\_\_hip\_bfloat16}, not the legacy \texttt{<hip/hip\_bfloat16.h>}).
(2) Q8\_0 repack-for-both-paths: serve decode \emph{and} prefill from the flat repacked int8$+$f16-scale arrays so the proven $739$-tok/s-on-4B coalesced-weight win transfers to 8B without doubling VRAM.
(3) Native Q4\_K/Q6\_K/IQ prefill through \texttt{qmmq\_core}.
(4) Relax the Q4\_K $k\%256$ construction guard so FFN widths like $17408$ stop falling back to dense f16.

\paragraph{ASPIRATIONAL.}
(a) Surpassing llama on this APU needs the ROCm-PagedAttention keystone (now built) plus split-KV flash decode for long-context occupancy and speculative decode/MTP.
(b) Long-context decode (split-KV / flash-split-k to recover the d512$\to$d1024 falloff).
(c) Backend portability via CubeCL: spiked, verdict NOT-YET---CubeCL's HIP compiler does not yet emit the int8 \texttt{iu8} intrinsic (it panics on int8) though the hardware has it; viable today for f16/bf16 and Vulkan-int8 only. Keep the hand-tuned HIP int8 path until a tuned CubeCL kernel reaches $\geq 60$--$70\%$ of it across backends.

\section{Method and Lessons}

The campaign reduces to a handful of transferable lessons, each earned the hard way. (1) \textbf{Measure the real kernel, not a proxy:} the $2.37\times$ decode lever was found only because we profiled the kernel actually wired into the Rust decode path and discovered it was a different, slower kernel than the look-alike microbench. (2) \textbf{Roofline first; let the ceiling tell you when to stop:} the matvec is at $97\%$ of the wall, so it is done; llama hitting the same $25$ tok/s confirmed $25$ is physics. (3) \textbf{Copy llama's actual config; don't blind-tune:} reading \texttt{mmq.cuh} line by line beat constant-sweeping---and told us which of llama's choices ($8$-warp, stream-K) to reject on this chip, with ISA evidence. (4) \textbf{Build the per-kernel numeric oracle before the kernel:} the \texttt{iu8} signedness NaN was undebuggable in a 36-layer forward; a tiny GPU-versus-exact-CPU test isolated it in one shot. (5) \textbf{Re-measure the bottleneck after every major change:} the entire hipGraph project was correct, exhaustively debugged, and useless because the fusion work had already removed the launches it was built to collapse. (6) \textbf{On a shared-memory APU, bandwidth is the budget:} every resident dense-f16 copy ``optimization'' backfired (decode $\to 5$ tok/s; the 8B repack spill $\to 58$ tok/s).

The bottom line is deliberately unembellished: on a consumer RDNA3.5 APU with no vendor matmul library, a measurement-driven path took native-ROCm inference from the slowest backend in the project to effective decode parity ($0.98\times$) and $\approx 0.86\times$ prefill against \texttt{llama.cpp} on the same silicon, while a unified 1-bit-to-full quant core kept every format bit-exact. We do not beat llama here; the remaining prefill gap is one well-scoped multi-wave WMMA flash-attention kernel, and surpassing llama is a keystone-project (ROCm PagedAttention, now built) away---both stated as engineering, not marketing.

\begin{thebibliography}{10}

\bibitem{ggml}
G. Gerganov et al. \textit{ggml: Tensor library for machine learning}. 2023.

\bibitem{llamacpp}
G. Gerganov et al. \textit{llama.cpp: LLM inference in C/C++}. 2023.

\bibitem{qwen3}
Qwen Team. \textit{Qwen3 Technical Report}. Alibaba, 2025.

\bibitem{williams2009}
S. Williams, A. Waterman, and D. Patterson. Roofline: An Insightful Visual Performance Model for Multicore Architectures. \textit{Communications of the ACM}, 52(4):65--76, 2009.

\bibitem{kwon2023}
W. Kwon et al. Efficient Memory Management for Large Language Model Serving with PagedAttention. \textit{SOSP}, 2023.

\bibitem{dao2022}
T. Dao et al. FlashAttention: Fast and Memory-Efficient Exact Attention with IO-Awareness. \textit{NeurIPS}, 2022.

\bibitem{rocm2024}
AMD. \textit{ROCm: Open Software Platform for GPU Compute}. Version 7.1, 2026.

\end{thebibliography}

\end{document}
